\section{Protocolo de simulación y registro de evidencias}
\label{sec:protocolo-simulacion}

La simulación permite contrastar las preferencias de los interesados ante funcionalidades del sistema de matrícula y calificaciones de una unidad educativa. El instrumento Kano establece una pregunta funcional (presencia de la característica) y una disfuncional (ausencia), usando cinco opciones comunes \cite{pohl2025,wiegers2013}. El protocolo describe una aplicación prevista; no constituye un informe de encuestas ya ejecutadas.

\subsection{Características de evaluación Kano}

Se seleccionaron ocho características relacionadas con los veinte requisitos candidatos. La correspondencia entre el código Kano y el código de requisito se conserva en \path{datos/caracteristicas_kano.csv}.

\begingroup
\small
\begin{longtable}{@{}p{0.15\linewidth}p{0.18\linewidth}p{0.61\linewidth}@{}}
\caption{Correspondencia entre características Kano y requisitos escolares}\label{tab:kano-caracteristicas}\\
\toprule
\textbf{Código} & \textbf{Requisito} & \textbf{Característica} \\
\midrule
K01 & RF-04 & Matrícula en línea \\
K02 & RF-06 & Comprobante digital de matrícula \\
K03 & RF-08 & Registro de calificaciones \\
K04 & RF-10 & Consulta personal de calificaciones \\
K05 & RF-11 & Boletín académico digital \\
K06 & RNF-01 & Rapidez de matrícula \\
K07 & RNF-04 & Facilidad de uso \\
K08 & RNF-05 & Compatibilidad con móvil \\
\bottomrule
\end{longtable}
\endgroup

\subsection{Procedimiento de aplicación}

Se prevé aplicar el instrumento a cinco personas adultas distintas, vinculadas con actividades de matrícula o seguimiento escolar, identificadas mediante los códigos P01 a P05. Los roles pertinentes al dominio incluyen secretaría, docentes, representantes, autoridades y estudiantes adultos. Para cada característica, el participante responderá la versión funcional y la disfuncional de una pregunta. El instrumento incluye dieciséis preguntas y se encuentra en \path{datos/instrumento_kano.csv}.

Las respuestas se registrarán en el archivo CSV de Kano. Cada fila identificará el código del participante, la característica evaluada y sus dos respuestas. La estructura de los campos está descrita en \path{datos/ESQUEMA.md}. Los cinco participantes deben aportar ocho pares válidos cada uno, para un mínimo de cuarenta registros. El análisis posterior utilizará la matriz de Kano y distinguirá las categorías básica, unidimensional, atractiva, indiferente, inversa y contradictoria. No se imputarán respuestas ausentes.

\subsection{Votación ponderada y negociación}

Se prevé una segunda técnica con tres personas adultas diferentes de las participantes en Kano, identificadas como P06, P07 y P08. Cada participante distribuirá exactamente cien puntos entre los veinte requisitos del catálogo; sus votos se registrarán en \path{datos/votos.csv}. La clasificación MoSCoW será discutida por el equipo y documentada en \path{datos/moscow.csv}, incluyendo para cada requisito una prioridad, una estimación positiva de esfuerzo y una justificación. El esfuerzo acumulado del grupo Must no podrá superar el sesenta por ciento del esfuerzo total acordado. La comparación y las decisiones finales no deben atribuirse a participantes sin datos recopilados.

\subsection{Consentimiento, custodia y manifiesto}

La invitación a participar debe explicar el propósito académico, el carácter voluntario del ejercicio, el tratamiento de respuestas anonimizadas y la custodia privada de los respaldos. No se publicarán formularios firmados, audios, videos, nombres, teléfonos, correos personales ni documentos de identificación. El equipo conservará los originales en una carpeta privada accesible a las personas autorizadas y registrará sus metadatos en \path{evidencias/MANIFIESTO.md}, de acuerdo con prácticas prudentes de protección de información personal \cite{lopdp2021,acm2018}.

Cada fila del manifiesto se incorporará únicamente al disponer del archivo original. Registrará código del participante, rol en el dominio educativo, técnica, fecha de recolección, nombre del archivo privado y huella SHA-256 obtenida de sus bytes. El manifiesto público no contiene entradas de participantes hasta que existan evidencias efectivamente obtenidas.
